% GENERATED FILE. Edit canonical lessons, then run npm run generate:books.
% canonical-source-hash: fnv1a64:71fac6a4a66b9be5
% canonical-lessons: UR-C505-raddi, UR-C505-jala, UR-C505-atari, UR-C505-pharna, UR-C505-rakhvana

\chapter{Up in the Loft}
\label{ch:ur-up-in-the-loft}
\hlchaptermodality{505}

\begin{chapteropening}
\textbf{By the end of this chapter:} \emph{I can talk about clearing out old things.}

Complete the last lesson of chapter 505: I can talk about clearing out old things.
\end{chapteropening}

\section[\texorpdfstring{\ur{ردی} (raddī)}{ردی (raddī)}]{\ur{ردی} (raddī) — waste paper; scrap}

\label{lesson:UR-C505-raddi}



\begin{quote}
Before the new one: say the Urdu for a peg (on a wall), then the Urdu for a socket (for a plug).
\end{quote}

\subsection*{You'll want to know: \ur{ردی}}
\textbf{\ur{ردی}} — \emph{raddī} — \textquotedblleft{}waste paper; scrap\textquotedblright{}.

\begin{quote}
\ur{ہم} \ur{پرانے} \ur{اخبار} \ur{ردی} \ur{میں} \ur{بیچتے} \ur{ہیں۔}
\end{quote}

\begin{quote}
\emph{ham purāne akhbār raddī meṅ bechte haiṅ.}
\end{quote}

\begin{quote}
We sell our old newspapers as waste paper.
\end{quote}

Feminine.

\begin{cousinweb}
\textbf{\ur{ردی}} (\emph{raddī}) is usually traced to Persian \emph{raddī}, \textquotedblleft{}rejected, worthless\textquotedblright{}.
\end{cousinweb}

\begin{grammarlens}[title={What you've built}]
One more word for describing what you lived and felt.
\end{grammarlens}

\subsection*{Your turn}
\begin{itemize}
\raggedright
  \item \emph{Say it:} \emph{raddī}
  \item \emph{Say it:} \emph{raddī}, once more
  \item \emph{Say it:} \emph{ham purāne akhbār raddī meṅ bechte haiṅ.} — We sell our old newspapers as waste paper.
  \item \emph{Recall:} say the Urdu for a peg (on a wall), then the Urdu for a socket (for a plug), then say \emph{raddī} again
\end{itemize}

\begin{tcolorbox}[breakable,colback=teal!4,colframe=teal!35!black,title={Before you move on}]
What does \ur{ردی} mean? (\textquotedblleft{}Waste paper; scrap\textquotedblright{}.) Say it once more.
\end{tcolorbox}
\section[\texorpdfstring{\ur{جالا} (jālā)}{جالا (jālā)}]{\ur{جالا} (jālā) — a cobweb}

\label{lesson:UR-C505-jala}



\begin{quote}
Before the new one: say the Urdu for a socket (for a plug), then the Urdu for waste paper.
\end{quote}

\subsection*{You'll want to know: \ur{جالا}}
\textbf{\ur{جالا}} — \emph{jālā} — \textquotedblleft{}a cobweb\textquotedblright{}.

\begin{quote}
\ur{چھت} \ur{کے} \ur{کونے} \ur{میں} \ur{جالا} \ur{ہے۔}
\end{quote}

\begin{quote}
\emph{chhat ke kone meṅ jālā hai.}
\end{quote}

\begin{quote}
There is a cobweb in the corner of the ceiling.
\end{quote}

Masculine.

\begin{cousinweb}
\textbf{\ur{جالا}} (\emph{jālā}) is inherited: CDIAL lists the Hindi cousin \emph{jālā}, \textquotedblleft{}a net, a web\textquotedblright{}, under Sanskrit \emph{jāla}, \textquotedblleft{}a net\textquotedblright{} (CDIAL 5213).
\end{cousinweb}

\begin{grammarlens}[title={What you've built}]
One more word for describing what you lived and felt.
\end{grammarlens}

\subsection*{Your turn}
\begin{itemize}
\raggedright
  \item \emph{Say it:} \emph{jālā}
  \item \emph{Say it:} \emph{jālā}, once more
  \item \emph{Say it:} \emph{chhat ke kone meṅ jālā hai.} — There is a cobweb in the corner of the ceiling.
  \item \emph{Recall:} say the Urdu for a socket (for a plug), then the Urdu for waste paper, then say \emph{jālā} again
\end{itemize}

\begin{tcolorbox}[breakable,colback=teal!4,colframe=teal!35!black,title={Before you move on}]
What does \ur{جالا} mean? (\textquotedblleft{}A cobweb\textquotedblright{}.) Say it once more.
\end{tcolorbox}
\section[\texorpdfstring{\ur{اٹاری} (aṭārī)}{اٹاری (aṭārī)}]{\ur{اٹاری} (aṭārī) — a loft; an attic}

\label{lesson:UR-C505-atari}



\begin{quote}
Before the new one: say the Urdu for waste paper, then the Urdu for a cobweb.
\end{quote}

\subsection*{You'll want to know: \ur{اٹاری}}
\textbf{\ur{اٹاری}} — \emph{aṭārī} — \textquotedblleft{}a loft; an attic\textquotedblright{}.

\begin{quote}
\ur{ہمارا} \ur{پرانا} \ur{سامان} \ur{اٹاری} \ur{میں} \ur{ہے۔}
\end{quote}

\begin{quote}
\emph{hamārā purānā sāmān aṭārī meṅ hai.}
\end{quote}

\begin{quote}
Our old things are up in the loft.
\end{quote}

Feminine.

\begin{cousinweb}
\textbf{\ur{اٹاری}} (\emph{aṭārī}) is inherited: CDIAL lists the Hindi cousin \emph{aṭārī}, \textquotedblleft{}a thatched upper building\textquotedblright{}, under Sanskrit \emph{aṭṭāla}, \textquotedblleft{}a watch-tower\textquotedblright{} (CDIAL 185).
\end{cousinweb}

\begin{grammarlens}[title={What you've built}]
One more word for describing what you lived and felt.
\end{grammarlens}

\subsection*{Your turn}
\begin{itemize}
\raggedright
  \item \emph{Say it:} \emph{aṭārī}
  \item \emph{Say it:} \emph{aṭārī}, once more
  \item \emph{Say it:} \emph{hamārā purānā sāmān aṭārī meṅ hai.} — Our old things are up in the loft.
  \item \emph{Recall:} say the Urdu for waste paper, then the Urdu for a cobweb, then say \emph{aṭārī} again
\end{itemize}

\begin{tcolorbox}[breakable,colback=teal!4,colframe=teal!35!black,title={Before you move on}]
What does \ur{اٹاری} mean? (\textquotedblleft{}A loft; an attic\textquotedblright{}.) Say it once more.
\end{tcolorbox}
\section[\texorpdfstring{\ur{پھاڑنا} (phāṛnā)}{پھاڑنا (phāṛnā)}]{\ur{پھاڑنا} (phāṛnā) — to tear (something)}

\label{lesson:UR-C505-pharna}



\begin{quote}
Before the new one: say the Urdu for a cobweb, then the Urdu for a loft.
\end{quote}

\subsection*{You'll want to know: \ur{پھاڑنا}}
\textbf{\ur{پھاڑنا}} — \emph{phāṛnā} — \textquotedblleft{}to tear (something)\textquotedblright{}.

\begin{quote}
\ur{بچے} \ur{نے} \ur{پرانا} \ur{اخبار} \ur{پھاڑا۔}
\end{quote}

\begin{quote}
\emph{bachche ne purānā akhbār phāṛā.}
\end{quote}

\begin{quote}
The child tore up the old newspaper.
\end{quote}

Past, with \ur{نے} (he or she …): \ur{اس} \ur{نے} \ur{خط} \ur{پھاڑا}, \emph{us ne khat phāṛā}.

\begin{cousinweb}
\textbf{\ur{پھاڑنا}} (\emph{phāṛnā}) is inherited: CDIAL lists the Hindi cousin \emph{phāṛnā}, \textquotedblleft{}to split\textquotedblright{}, under Sanskrit \emph{sphāṭayati}, \textquotedblleft{}causes to split\textquotedblright{} (CDIAL 13825).
\end{cousinweb}

\begin{grammarlens}[title={What you've built}]
One more word for describing what you lived and felt.
\end{grammarlens}

\subsection*{Your turn}
\begin{itemize}
\raggedright
  \item \emph{Say it:} \emph{phāṛnā}
  \item \emph{Say it:} \emph{phāṛnā}, once more
  \item \emph{Say it:} \emph{bachche ne purānā akhbār phāṛā.} — The child tore up the old newspaper.
  \item \emph{Recall:} say the Urdu for a cobweb, then the Urdu for a loft, then say \emph{phāṛnā} again
\end{itemize}

\begin{tcolorbox}[breakable,colback=teal!4,colframe=teal!35!black,title={Before you move on}]
What does \ur{پھاڑنا} mean? (\textquotedblleft{}To tear (something)\textquotedblright{}.) Say it once more.
\end{tcolorbox}
\section[\texorpdfstring{\ur{رکھوانا} (rakhvānā)}{رکھوانا (rakhvānā)}]{\ur{رکھوانا} (rakhvānā) — to have (something) put away}

\label{lesson:UR-C505-rakhvana}



\begin{quote}
Before the new one: say the Urdu for a loft, then the Urdu for to tear (something).
\end{quote}

\subsection*{You'll want to know: \ur{رکھوانا}}
\textbf{\ur{رکھوانا}} — \emph{rakhvānā} — \textquotedblleft{}to have (something) put away\textquotedblright{}.

\begin{quote}
\ur{امی} \ur{نے} \ur{ردی} \ur{کے} \ur{اخبار} \ur{باہر} \ur{رکھوائے۔}
\end{quote}

\begin{quote}
\emph{ammī ne raddī ke akhbār bāhar rakhvāe.}
\end{quote}

\begin{quote}
Mum had the old newspapers put outside.
\end{quote}

Past, with \ur{نے} (he or she …): \ur{اس} \ur{نے} \ur{سامان} \ur{گاڑی} \ur{میں} \ur{رکھوایا}, \emph{us ne sāmān gāṛī meṅ rakhvāyā}.

\begin{cousinweb}
\textbf{\ur{رکھوانا}} (\emph{rakhvānā}) is the causative of \emph{rakhnā}, \textquotedblleft{}to put, to keep\textquotedblright{}, which CDIAL lists under Sanskrit \emph{rakṣati}, \textquotedblleft{}guards\textquotedblright{} (CDIAL 10547): to have someone put a thing away.
\end{cousinweb}

\begin{grammarlens}[title={What you've built}]
One more word for describing what you lived and felt.
\end{grammarlens}

\subsection*{Your turn}
\begin{itemize}
\raggedright
  \item \emph{Say it:} \emph{rakhvānā}
  \item \emph{Say it:} \emph{rakhvānā}, once more
  \item \emph{Say it:} \emph{ammī ne raddī ke akhbār bāhar rakhvāe.} — Mum had the old newspapers put outside.
  \item \emph{Recall:} say the Urdu for a loft, then the Urdu for to tear (something), then say \emph{rakhvānā} again
\end{itemize}

\begin{tcolorbox}[breakable,colback=teal!4,colframe=teal!35!black,title={Before you move on}]
What does \ur{رکھوانا} mean? (\textquotedblleft{}To have (something) put away\textquotedblright{}.) Say it once more.
\end{tcolorbox}
